\chapter{Background and Related Work}
\label{ch:background}

\section{Chunking in retrieval-augmented generation}
\label{sec:rag_chunking}

Standard large language models (LLMs) struggle with knowledge-intensive tasks because their
internal knowledge is limited by a static training cut-off and their output lacks
verifiable sources~\cite{lewis2020rag}. Retrieval-augmented generation (RAG) resolves these
limitations by coupling pre-trained parametric generators (such as BART) with
non-parametric memory in the form of an external dense vector index~\cite{lewis2020rag}. In
a RAG pipeline, raw source documents are first split into retrievable text units, the
so-called \emph{chunks}~\cite{smigielski2026chunking}, which are then converted into dense
vector representations by a neural encoder and finally searched through maximum inner
product search (MIPS) or approximate nearest-neighbour algorithms~\cite{lewis2020rag}.

This segmentation is a pre-processing step that determines the quality of the subsequent
vector index~\cite{smigielski2026chunking}. Downstream components, such as cross-encoder
rerankers, have to rely on the initial segmentation; errors made during this step cannot be
repaired later, and the subsequent components are forced to process chunks that are either
truncated or contextually degraded~\cite{duarte2024lumberchunker}.

Chunking can fail in two ways. Chunks that are too small or cut mid-sentence lack
sufficient factual context, which leads to retrieval failures and a higher risk of
hallucinations by the LLM. Conversely, chunks that are too large contain irrelevant
content, which dilutes the vector embedding and creates a semantic misalignment with the
query vector during dense retrieval~\cite{smigielski2026chunking}.

Despite its decisive role in a RAG pipeline, chunking has historically been treated as an
arbitrary pre-processing step done through fixed character cuts rather than as an
architectural variable~\cite{qu-etal-2025-semantic}. For comparative studies, recent
literature emphasises the need to isolate chunking from the rest of the RAG pipeline in
order to guarantee consistency across the compared splitting
variants~\cite{smigielski2026chunking}.

\section{Rule-based and embedding-based strategies}
\label{sec:rule_based}

\begin{figure}[htbp]
  \centering
  \begin{tikzpicture}[x=0.9cm,y=1cm,
      sent/.style={draw=black!60, line width=0.3pt},
      cut/.style={line width=1.6pt, black},
      badcut/.style={line width=1.6pt, red!75!black},
      lbl/.style={anchor=west, font=\small},
      note/.style={font=\scriptsize, text=black!70}]
    \colorlet{topicA}{blue!15}
    \colorlet{topicB}{orange!30}
    % one row of eight sentences; #1 = y position of the row
    \newcommand{\sentstrip}[1]{%
      \foreach \a/\b/\t/\n in {0/1.6/A/1, 1.6/3.8/A/2, 3.8/5.2/A/3, 5.2/7.0/A/4,
                               7.0/9.0/B/5, 9.0/10.5/B/6, 10.5/12.6/B/7, 12.6/14.0/B/8}{%
        \draw[sent, fill=topic\t] (\a,#1) rectangle (\b,#1+0.5);
        \node[font=\scriptsize] at ({(\a+\b)/2},#1+0.25) {$s_{\n}$};}}
    \newcommand{\cutat}[3][cut]{\draw[#1] (#2,#3-0.15) -- (#2,#3+0.65);}

    % document
    \node[lbl] at (0,11.35) {Document: eight sentences, two topics, three paragraphs};
    \sentstrip{10.5}
    \foreach \x in {3.8,10.5} \node[font=\small] at (\x,10.25) {\P};
    \draw[decorate, decoration={brace, mirror, amplitude=4pt}] (0,10.0) -- (7.0,10.0)
      node[midway, below=4pt, note] {topic A};
    \draw[decorate, decoration={brace, mirror, amplitude=4pt}] (7.0,10.0) -- (14.0,10.0)
      node[midway, below=4pt, note] {topic B};

    % (a) fixed-size
    \node[lbl] at (0,8.55) {(a) Fixed-size: a cut every $L$ characters};
    \sentstrip{7.7}
    \foreach \x in {4.0,8.0,12.0} \cutat[badcut]{\x}{7.7}
    \draw[<->, note] (0,7.35) -- (4.0,7.35) node[midway, below=-1pt] {$L$};

    % (b) recursive
    \node[lbl] at (0,6.15) {(b) Recursive: paragraph break $\rightarrow$ line break
      $\rightarrow$ space, as long as a piece is longer than $L$};
    \sentstrip{5.3}
    \cutat{3.8}{5.3} \cutat{10.5}{5.3} \cutat[badcut]{7.8}{5.3}
    \node[note, anchor=north] at (7.8,5.15) {no \P\ within $L$: cut at a space};

    % (c) semantic
    \node[lbl] at (0,3.75) {(c) Semantic: cut where the embedding distance of neighbours
      exceeds a threshold};
    \sentstrip{2.9}
    \foreach \x/\h in {1.6/0.15, 3.8/0.5, 5.2/0.2, 7.0/0.8, 9.0/0.2, 10.5/0.3, 12.6/0.15}
      \fill[black!45] (\x-0.12,1.85) rectangle (\x+0.12,1.85+\h);
    \draw[dashed, red!75!black] (0,2.25) -- (14.0,2.25)
      node[right, note, text=red!75!black] {threshold};
    \node[note, anchor=east] at (0,2.05) {distance};
    \cutat{3.8}{2.9} \cutat{7.0}{2.9}

    % (d) LLM-based
    \node[lbl] at (0,1.05) {(d) LLM-based: the model is asked whether the next sentences
      continue the topic};
    \sentstrip{0.2}
    \foreach \x in {1.6,3.8,5.2,9.0,10.5,12.6}
      \node[note, text=green!45!black] at (\x,-0.1) {yes};
    \node[note, text=red!75!black, font=\scriptsize\bfseries] at (7.0,-0.1) {no};
    \cutat{7.0}{0.2}

    % legend
    \draw[cut] (0.1,-0.75) -- (0.1,-0.45);
    \node[lbl, font=\scriptsize] at (0.2,-0.6) {boundary between sentences};
    \draw[badcut] (5.2,-0.75) -- (5.2,-0.45);
    \node[lbl, font=\scriptsize] at (5.3,-0.6) {boundary inside a sentence};
    \node[lbl, font=\scriptsize] at (10.2,-0.6) {\P\ paragraph break};
  \end{tikzpicture}
  \caption{The chunking strategies of \cref{sec:rule_based,sec:llm_chunking}, applied to the
  same schematic document. Fixed-size splitting (a) ignores the text and cuts inside
  sentences. Recursive splitting (b) prefers paragraph breaks but falls back to spaces
  inside a long paragraph. Semantic chunking (c) cuts where consecutive sentence
  embeddings differ strongly, which can also produce boundaries within a topic ($s_2|s_3$).
  LLM-based chunking (d) asks a language model at each step and cuts only at the topic
  change.}
  \label{fig:strategies}
\end{figure}

\Cref{fig:strategies} illustrates the strategies discussed in this chapter on one schematic
document. \emph{Fixed-size chunking} (\cref{fig:strategies}a) splits a source document
into uniform segments based on a predetermined character or token
limit~\cite{qu-etal-2025-semantic}. Although the
procedure is computationally simple and deterministic, it disregards syntax and thematic
organisation~\cite{smigielski2026chunking}. This leads to chunks that are cut mid-topic and
thus to a potential degradation in retrieval quality~\cite{qu-etal-2025-semantic}. To
mitigate these effects, RAG implementations let consecutive fixed-size segments overlap in
a sliding window, typically by \SIrange{10}{20}{\percent} of the chunk
length~\cite{smigielski2026chunking}. However, a recent empirical evaluation found that
overlap provides no measurable improvement in answer quality and only increases the
indexed text and the memory overhead~\cite{bennani2026chunking}.

In contrast, \emph{recursive character splitting} (\cref{fig:strategies}b) traverses the document through an
ordered hierarchy of typographical separators and thereby preserves its structure where
possible~\cite{duarte2024lumberchunker}. Typically, such splitters cascade through
paragraph breaks, line breaks, sentence punctuation and finally spaces: when a segment
still exceeds the size limit, the algorithm moves one step down the hierarchy to find a
split point~\cite{duarte2024lumberchunker}.

Embedding-based strategies differ from these typographical heuristics in that they detect
boundaries semantically. \emph{Semantic chunking} (\cref{fig:strategies}c) converts individual sentences or passages
into dense vector embeddings, numerical representations in a high-dimensional space,
using models such as Sentence-BERT (SBERT)~\cite{reimers2019sbert}. The cosine distance
between consecutive segments then determines whether a thematic shift has occurred. More
advanced approaches such as Max--Min semantic chunking compare the similarity of each
incoming sentence to the current chunk with the minimum similarity within that chunk, so
that the text is split without a fixed window size~\cite{kiss2025maxmin,
smigielski2026chunking}. However, the benefit of semantic chunking depends strongly on the
data. On artificially stitched datasets with high topic diversity, it preserves topic
boundaries better than recursive algorithms, but on real-world documents, which already
provide natural cues for topic boundaries, the retrieval accuracy does not differ
substantially while the computational cost is higher~\cite{qu-etal-2025-semantic}. The
additional cost, latency and memory needed to compute sentence-level embeddings are
therefore rarely justified for standard structured text~\cite{qu-etal-2025-semantic}.

\section{LLM-based chunking strategies}
\label{sec:llm_chunking}

LLM-based segmentation strategies replace distance thresholds with generative models that
identify semantic shifts in a document (\cref{fig:strategies}d)~\cite{smigielski2026chunking}. Earlier work in this
field includes LumberChunker, which divides a text into paragraphs and assigns them
sequential IDs starting at 1. Consecutive paragraphs are collected into a group until a
target token count is exceeded, and the group is then submitted to an LLM with a prompt
that asks for the ID of the paragraph at which the narrative or topical focus
changes~\cite{duarte2024lumberchunker}.

On a benchmark of \num{3000} question--answer pairs from 100 narrative books of Project
Gutenberg, LumberChunker reached a Recall@20, the share of relevant passages retrieved
among the top 20 chunks, of \SI{77.92}{\percent}, outperforming the recursive chunking
baseline with \SI{74.35}{\percent}~\cite{duarte2024lumberchunker}.

However, these improvements in retrieval accuracy come with high operational and
computational costs~\cite{smigielski2026chunking}. LLM-based chunking is generally the most
expensive form of pre-processing. Because each prompt depends on the boundary returned for the previous
passage, the queries cannot be parallelised~\cite{duarte2024lumberchunker}. The processing
time therefore grows with the size of the document~\cite{duarte2024lumberchunker}.

\section{Shortcomings in current literature and methodological commitments}
\label{sec:commitments}

Existing literature on document segmentation for RAG shows substantial methodological
limitations. First, comparative evaluations frequently test chunking strategies without
calibrating the baselines to the average chunk length of the LLM
splitters~\cite{qu-etal-2025-semantic}. For example, in the evaluation of LumberChunker,
the recursive baselines averaged 399 tokens per chunk, while the LLM chunker produced a
mean of 334 tokens~\cite{duarte2024lumberchunker}. Without normalising the baseline
lengths, the quality of the boundary placement is confounded with chunk-size effects,
since chunk size directly affects retrieval recall and the density of relevant
context~\cite{qu-etal-2025-semantic}. Second, benchmark studies also aggregate retrieval
and generation metrics without isolating the structural layout of a document as an
independent variable~\cite{qu-etal-2025-semantic, smigielski2026chunking}. While LLM
chunkers demonstrate advantages on long narrative prose~\cite{duarte2024lumberchunker},
their benefits diminish on real-world, structured text, where natural typographical
conventions already provide a good indication of boundary
placement~\cite{qu-etal-2025-semantic}. To address these gaps and to isolate the effect of
chunking, this thesis follows four methodological commitments:
\begin{enumerate}
    \item \textbf{Length-matched baselines.} All baselines are tuned to the mean chunk
    length produced by the LLM chunker on each document (\cref{sec:design}). This removes
    the chunk-size bias and isolates the quality of the boundary
    placement~\cite{qu-etal-2025-semantic}.

    \item \textbf{Document structure as an independent variable.} The evaluated documents
    belong to three structure classes (explicit, implicit and unstructured,
    \cref{sec:datasets}), in order to evaluate under which structural conditions LLM-based
    chunking justifies its computational overhead~\cite{qu-etal-2025-semantic,
    smigielski2026chunking}.

    \item \textbf{Local execution and determinism.} To make the project usable for
    everybody and to work within the constraints of the available hardware, no commercial
    API is used; all segmentation runs use an open-weight local model with deterministic
    decoding ($T = 0$ and a fixed seed)~\cite{smigielski2026chunking}.

    \item \textbf{Pre-specified analysis.} The metrics, the paired significance test and
    the correction for multiple testing were fixed in a written test protocol before the
    final chunking runs; every decision taken after results were known is reported as such
    (\cref{sec:threats}).
\end{enumerate}
